%
% spec-writeset-admission.tex
%
% Write-Set Admission: A Z Specification
% A formal model of path canonicalization and write-set admission as
% implemented in internal/mission/conflict.go and consumed by
% internal/hook/pretooluse.go, current as of the commit that
% introduced this document (branch feat/z-spec-state-machines).
%
% Compile: pdflatex spec-writeset-admission.tex
%
% Must be .tex, not .md, for the same reason given in
% spec-mission-lifecycle.tex's own header comment: this repo's
% /z-spec:test tooling globs for .tex specifically, and probcli's
% Z-to-B translation requires the extension to run at all.
%
\documentclass[a4paper,10pt,fleqn]{article}
\usepackage[margin=1in]{geometry}
\usepackage{fuzz}
\usepackage{booktabs}
\usepackage{tabularx}
\usepackage[colorlinks=true,linkcolor=blue,citecolor=blue,urlcolor=blue]{hyperref}

\newcommand{\fpath}[1]{\texttt{#1}}
\newcommand{\field}[1]{\texttt{#1}}

\begin{document}

\title{Write-Set Admission: A Z Specification}
\author{Mike S, for Punt Labs}
\date{September 2026}
\hypersetup{
  pdftitle={Write-Set Admission: A Z Specification},
  pdfauthor={Punt Labs},
  pdfsubject={Z Specification},
  pdfcreator={fuzz/probcli}
}
\maketitle
\tableofcontents
\newpage

% ===================================================================
\section{Scope}
\label{sec:scope}
% ===================================================================

This specification models two things that live together in
\fpath{internal/mission/conflict.go} because a mistake in either one
opens the same wound --- a worker writing where its contract did not
authorize it:

\begin{enumerate}
  \item \emph{Path canonicalization} --- \field{splitSegments}
    (conflict.go:475--497) and \field{CanonicalPath}
    (conflict.go:518--524), which turn a raw path string into the
    segment list every containment and overlap decision is computed
    from; and
  \item \emph{Write-set admission} --- \field{findWriteSetConflicts}
    (conflict.go:70--120) and the six-rule dispatch underneath it,
    \field{anyEntryConflicts}/\field{entryPairConflicts}
    (conflict.go:154--217), which \field{Store.Create} consults to
    refuse a new mission contract whose declared write territory
    overlaps an existing \emph{open} mission's.
\end{enumerate}

The second of these is explicitly the property
spec-mission-lifecycle.tex's own Non-goals section named and set
aside: ``\emph{Cross-mission write-set conflicts --- ...\ This is a
Create-time admission check over the whole store, not a property of
one mission's lifecycle.}'' (spec-mission-lifecycle.tex
\S~Non-goals). This document picks up exactly that thread.

Two real defects motivate the Validation section
(\S\ref{sec:validation}): \field{ethos-vaib}, in which
\field{splitSegments} discards the leading empty segment an absolute
path produces, so \field{CanonicalPath("/docs/x.md")} and
\field{CanonicalPath("docs/x.md")} are the same string; and
\field{ethos-8ady}, in which mission lint's write-set coverage
heuristic (\fpath{internal/mission/lint.go}:191--197) treats a glob
entry no differently from a literal one --- its trailing-slash guard
(line~193) and its own \field{HasPrefix} comparison both operate on
\field{CanonicalPath}'s joined string, blind to what any segment
inside it means --- so a file genuinely covered by \texttt{docs/**}
draws a spurious ``not in write\_set'' warning.

\subsection{Non-goals}
\label{sec:nongoals}

\begin{itemize}
  \item \textbf{Single-segment wildcards} (\texttt{*}, \texttt{?},
    and character classes within one segment,
    \field{segmentMatches}, conflict.go:433--442, via Go's
    \field{path.Match}) --- neither defect this document checks, nor
    the write-set overlap property, depends on matching \emph{within}
    a segment. Only the whole-segment \texttt{**} form
    (\field{PATSEG}'s \field{wildRestSeg}, \S\ref{sec:freetypes})
    is modelled.
  \item \textbf{A \texttt{**} at a non-trailing pattern position} ---
    the real \field{segmentsContain} (conflict.go:394--419) handles
    this via backtracking, but every entry in this codebase places
    \texttt{**} last when it appears at all. \field{ENTRY}
    (\S\ref{sec:state}) carries this as a standing invariant rather
    than as an unchecked assumption.
  \item \textbf{Case-insensitive filesystems} --- \field{pathsOverlap}'s
    own comment (conflict.go:236--238) already scopes this out:
    ``\emph{macOS case-insensitive HFS+ is a known divergence and not
    handled here}.''
  \item \textbf{A write\_set file entry whose own text happens to
    contain a glob metacharacter, compared on the \emph{file} side of
    a containment check} --- the real \field{pathContainedBy} always
    glob-interprets the \emph{entry} argument and never the
    \emph{file} argument (segmentMatches is only ever called with an
    entry segment as its first parameter). \field{PathContainedByBuggy}/
    \field{PathContainedByFixed} (\S\ref{sec:containment}) preserve
    this asymmetry. The admission model's \field{Store} invariant
    (\S\ref{sec:admission}), which compares two write-set entries
    against \emph{each other}, applies the same wildcard-aware
    position rule to both sides for economy; this is very slightly
    more permissive than \field{entryPairConflicts}'s literal-file-side
    reading in the one case where a file-shaped entry's own segment
    text is itself \texttt{**}, a combination no entry in this
    codebase exhibits.
  \item \textbf{Role overlap, generic \field{Store.Update}, pipelines} ---
    excluded for the same reasons spec-mission-lifecycle.tex's own
    Non-goals section gives; they do not bear on path semantics.
\end{itemize}

\subsection{Modelling choice: schemas, not free functions, over
sequences}
\label{sec:modellingchoice}

\field{CanonicalPath} and \field{PathContainedBy} are presented below
as schemas (\S\ref{sec:canonical}--\ref{sec:containment}) whose
predicate relates an input field to an output field, not as
free-standing \fuzz{} function declarations of the form
$f : \seq X \fun Y$. This is not a stylistic preference; it is
forced by a fact about how \texttt{probcli} translates a Z document
to B, confirmed by trial before this document's own admission model
was written. Declaring \emph{any} global constant --- via
\field{axdef} or \field{gendef} --- whose type mentions $\seq X$ in a
function or relation's domain causes \texttt{SETUP\_CONSTANTS} to
fail for the \emph{whole} machine, even when the operations actually
being checked never reference that constant: a single unused
$\seq X \rel \seq X$ declaration, added to an otherwise-working
document, turns a clean 128-state run into
\texttt{setup\_constants\_fails} with zero states explored. $\seq X$
is unbounded as a Z type regardless of how tightly an axiom
constrains the function's actual values, and B's constant-solving
phase must represent the constant's full \emph{declared} type before
it can apply any axiom at all.

Schemas do not have this problem: a schema is a predicate, not a
constant requiring existence-proof, and a bounded state component
typed $\seq X$ (guarded by its own $\# \leq \field{maxSegs}$
invariant, per \S\ref{sec:constants}) animates exactly as
spec-mission-lifecycle.tex's own \field{resultRounds} and
\field{reflected} components did. \field{CanonicalPath} is therefore
given as the schema relating a raw input to its canonical output;
reading it as a function means fixing the input and solving the
schema's predicate for the unique output it determines, which is
Z's standard schema-as-relation-graph reading, not a departure from
it.

% ===================================================================
\section{Basic Types}
\label{sec:basic}
% ===================================================================

\begin{zed}
SEG ::= segDocs | segFile
\end{zed}

\field{SEG} stands for a normalized path segment name --- the
counterpart of \field{HANDLE} in spec-mission-lifecycle.tex: two
concrete constructors are enough to exhibit every defect and every
admission rule this document checks, and probcli needs a
finitely-enumerated carrier to animate at all.

% ===================================================================
\section{Free Types}
\label{sec:freetypes}
% ===================================================================

\subsection{ZBOOL}

Carried over from spec-mission-lifecycle.tex \S~ZBOOL for the same
reason: \fuzz{} has no built-in boolean, and \field{BOOL}/\field{true}/
\field{false} collide with B keywords under ProB's translation.

\begin{zed}
ZBOOL ::= ztrue | zfalse
\end{zed}

\subsection{RAWSEG}

The result of splitting a raw path string on \texttt{"/"}
\emph{before} \field{splitSegments}'s own filtering
(conflict.go:475--497) removes anything. Go's
\field{strings.Split("/docs/x.md", "/")} yields
$\langle \mathtt{""}, \mathtt{"docs"}, \mathtt{"x.md"} \rangle$ --- the
leading empty token is exactly how an absolute path manifests
\emph{before} the filter that discards it runs. \field{rsEmpty}
models that token; \field{rsDot} models a literal \texttt{"."} token
(from \texttt{./foo} or \texttt{internal/./foo}, which
validate.go's own \field{TestValidate\_AcceptsSingleDotSegment}
deliberately accepts as legitimate path syntax, DES-031); every other
token is a real segment, \field{rsSeg}-wrapped.

\begin{zed}
RAWSEG ::= rsSeg \ldata SEG \rdata | rsEmpty | rsDot
\end{zed}

\subsection{PATSEG}

A write-set entry's segment, once split: either a literal name or the
\texttt{**} wildcard (\S\ref{sec:nongoals} excludes the single-segment
\texttt{*}/\texttt{?} forms).

\begin{zed}
PATSEG ::= litSeg \ldata SEG \rdata | wildRestSeg
\end{zed}

% ===================================================================
\section{Global Constants}
\label{sec:constants}
% ===================================================================

Bounding \field{SEG}'s use in a sequence to a small constant is what
lets probcli enumerate the state space at all
(\S\ref{sec:modellingchoice}); the shape of every predicate below is
independent of the concrete value, exactly as
spec-mission-lifecycle.tex's own \field{maxRounds}/\field{maxDelegations}
were chosen small purely for this reason.

\begin{axdef}
maxSegs : \nat
\where
maxSegs = 2
\end{axdef}

\field{maxOpen} bounds how many concurrently-open missions
\S\ref{sec:admission}'s \field{Store} tracks. There is no real
analogue --- the actual \field{Store} holds arbitrarily many open
missions --- exactly as spec-mission-lifecycle.tex's
\field{maxDelegations} had none.

\begin{axdef}
maxOpen : \nat
\where
maxOpen = 2
\end{axdef}

% ===================================================================
\section{Path Canonicalization}
\label{sec:canonical}
% ===================================================================

\field{unwrap} recovers the real segment a \field{RAWSEG} token
carries; it is undefined on \field{rsEmpty} and \field{rsDot} ---
those tokens carry no segment, by construction.

\begin{axdef}
unwrap : RAWSEG \pfun SEG
\where
\dom unwrap = \ran rsSeg \\
\forall s : SEG @ unwrap(rsSeg(s)) = s
\end{axdef}

\subsection{CanonicalPathBuggy}

Today's real function. \field{splitSegments} (conflict.go:475--497)
drops every \field{rsEmpty} and \field{rsDot} token via
\texttt{squash} over the filtered relation, exactly as
\field{CanonicalPath} (conflict.go:518--524) does when it joins the
result back into a string. Its component list is the point: this
schema has nowhere to put an absoluteness bit, because the real
function's return type, \texttt{[]string}, has nowhere to put one
either.

\begin{schema}{CanonicalPathBuggy}
raw : \seq RAWSEG \\
segs : \seq SEG
\where
\# raw \leq maxSegs \\
segs = squash(raw \comp unwrap)
\end{schema}

\subsection{CanonicalPathFixed}

\field{ethos-vaib}'s prescribed fix (bd ticket: ``\emph{Fix requires
teaching CanonicalPath about absoluteness}''), stated at exactly the
point the bug lives: the corrected function's codomain must be able
to express what the buggy one cannot. No function computing this bit
exists anywhere in the current code; a leading \field{rsEmpty} token
is exactly Go's own \field{filepath.IsAbs} test, transported through
the split.

\begin{schema}{CanonicalPathFixed}
raw : \seq RAWSEG \\
absolute : ZBOOL \\
segs : \seq SEG
\where
\# raw \leq maxSegs \\
absolute = (\IF \# raw > 0 \land raw(1) = rsEmpty \THEN ztrue \ELSE zfalse) \\
segs = squash(raw \comp unwrap)
\end{schema}

Both schemas agree on \field{segs} for every \field{raw} --- the fix
changes what is \emph{exposed}, not how segments are computed. That
equality is recorded as a theorem, \field{SegmentationUnchanged}, in
\S\ref{sec:theorems}.

% ===================================================================
\section{Write-Set Entries}
\label{sec:state}
% ===================================================================

\begin{schema}{ENTRY}
isDir, isEI : ZBOOL \\
pattern : \seq PATSEG \\
rootGlob : ZBOOL \\
truncLen : \nat \\
seg1, seg2 : PATSEG
\where
\# pattern \leq maxSegs \\
isEI = ztrue \implies isDir = ztrue \\
\forall i : 1 \upto \# pattern @ pattern(i) = wildRestSeg \implies i = \# pattern \\
rootGlob = (\IF \# pattern = 1 \land pattern(1) = wildRestSeg \THEN ztrue \ELSE zfalse) \\
truncLen = (\IF \# pattern > 0 \land pattern(\# pattern) = wildRestSeg
             \THEN \# pattern - 1 \ELSE \# pattern) \\
truncLen \geq 1 \implies seg1 = pattern(1) \\
truncLen \geq 2 \implies seg2 = pattern(2)
\end{schema}

Seven fields, three kinds:

\begin{itemize}
  \item \field{isDir}, \field{isEI}, \field{pattern} are the real
    data: \field{isDirEntry}'s trailing-slash test
    (conflict.go:219--226), the \field{ExtractInto} field's own
    entry-kind flag, and the entry's segments after
    \field{CanonicalPath}-style splitting.
  \item The predicate \field{isEI = ztrue \implies isDir = ztrue} is
    validate.go rule~17 (validate.go:79, 216--226): ``\emph{extract\_into
    entries must be directory-shaped}.'' The
    $(\lnot\field{isDir} \land \field{isEI})$ combination
    \field{entryPairConflicts}'s own comment calls ``unreachable'' and
    guards defensively anyway (conflict.go:183--185, 193--195); this
    document's invariant makes that unreachability a checked fact
    rather than a defensive comment.
  \item \field{rootGlob}, \field{truncLen}, \field{seg1}, \field{seg2}
    have no code counterpart. They are this document's
    animatability device (\S\ref{sec:modellingchoice}): every
    predicate that needs to compare two entries' patterns
    (\field{PathContainedByBuggy}/\field{Fixed},
    \S\ref{sec:containment}; \field{Store}'s invariant and
    \field{AdmitMission}'s guard, \S\ref{sec:admission}) reads them
    instead of re-deriving \field{overlapSegments}'s truncation
    (conflict.go:291--314) by hand at every use site. The invariant
    above bans \texttt{**} at any position but the last
    (\S\ref{sec:nongoals}), so \field{truncLen} is exactly the length
    of \field{pattern} with any trailing \texttt{**} removed, and
    \field{seg1}/\field{seg2} are its first two segments where they
    exist.
\end{itemize}

% ===================================================================
\section{Containment}
\label{sec:containment}
% ===================================================================

\subsection{PathContainedByBuggy}

Today's real \field{pathContainedBy} (conflict.go:360--379) and
\field{segmentsContain} (conflict.go:394--419), specialized to a
single trailing \texttt{**} (\S\ref{sec:nongoals}): the entry's
literal prefix (\field{truncLen}, \field{seg1}, \field{seg2}) must
match a prefix of the file's segments, and the file may be longer
--- ``\emph{the entry running out is success}'' (conflict.go:389).
Included via \field{CanonicalPathBuggy}, so the file side carries no
absoluteness information to check.

\begin{schema}{PathContainedByBuggy}
CanonicalPathBuggy \\
ENTRY \\
containedBy : ZBOOL
\where
(containedBy = ztrue \iff
  truncLen \leq \# segs \land
  (truncLen \geq 1 \implies seg1 = litSeg(segs(1))) \land
  (truncLen \geq 2 \implies seg2 = litSeg(segs(2))))
\end{schema}

\subsection{PathContainedByFixed}

The one-conjunct fix: a file path that is absolute cannot be
contained by a write-set entry, because every entry is repo-relative
by construction (validate.go:562--568 rejects an absolute
\field{write\_set}/\field{extract\_into} entry outright). This is
criterion (a) from the mission brief made concrete: canonicalization
preserving the absolute/relative distinction is only useful if
containment then \emph{consults} it.

\begin{schema}{PathContainedByFixed}
CanonicalPathFixed \\
ENTRY \\
containedBy : ZBOOL
\where
(containedBy = ztrue \iff
  absolute = zfalse \land
  truncLen \leq \# segs \land
  (truncLen \geq 1 \implies seg1 = litSeg(segs(1))) \land
  (truncLen \geq 2 \implies seg2 = litSeg(segs(2))))
\end{schema}

Neither schema compares \field{segs} to \field{pattern} as strings.
Both consult \field{seg1}/\field{seg2}/\field{truncLen}, which are
already glob-resolved by \field{ENTRY}'s own predicate --- this is
criterion (b): containment is path-semantic, not
pattern-string-literal. The contrast that fails this property ---
comparing a glob entry's segments as literal text --- is exhibited
as a reachable state in \S\ref{sec:validation-8ady}, not restated
here as a third schema, to avoid presenting the same
\field{ethos-8ady} contrast twice.

% ===================================================================
\section{Write-Set Admission}
\label{sec:admission}
% ===================================================================

Models \field{Store.Create}'s cross-mission conflict scan:
\field{findWriteSetConflicts} (conflict.go:70--120) dispatching
through \field{anyEntryConflicts}/\field{entryPairConflicts}
(conflict.go:154--217). \field{Overlaps}, the six-rule relation those
functions compute, is not given as a separate schema here for the
reason \S\ref{sec:modellingchoice} explains: it would need to relate
two \field{ENTRY} values as a free \fuzz{} relation, and
\field{ENTRY}'s own \field{pattern} field is $\seq PATSEG$-typed, so
the same \texttt{SETUP\_CONSTANTS} failure applies to a schema-typed
domain exactly as it does to a bare sequence type. It is instead
written out, once, inline in \field{Store}'s own invariant, and
\field{AdmitMission}'s guard is the calculated precondition of
preserving that invariant --- the same relationship
spec-mission-lifecycle.tex's \field{Abandon} has to
\field{AbandonSafety}.

\begin{schema}{Store}
open : 1 \upto maxOpen \pfun ENTRY
\where
\forall i, j : \dom open | i \neq j @
  \lnot ( \lnot ((open\ i).isEI = ztrue \land (open\ j).isEI = ztrue) \land
    ( ( (open\ i).isDir = (open\ j).isDir \land
        (open\ i).truncLen > 0 \land (open\ j).truncLen > 0 \land
        ( (open\ i).rootGlob = ztrue \lor (open\ j).rootGlob = ztrue
          \lor ( (open\ i).seg1 = (open\ j).seg1 \land
                ( (open\ i).truncLen = 1 \lor (open\ j).truncLen = 1
                  \lor (open\ i).seg2 = (open\ j).seg2 ) ) ) )
      \lor ( (open\ i).isDir = zfalse \land (open\ j).isDir = ztrue \land
        (open\ j).truncLen \leq (open\ i).truncLen \land
        ( (open\ j).truncLen \geq 1 \implies (open\ j).seg1 = (open\ i).seg1 ) \land
        ( (open\ j).truncLen \geq 2 \implies (open\ j).seg2 = (open\ i).seg2 ) )
      \lor ( (open\ j).isDir = zfalse \land (open\ i).isDir = ztrue \land
        (open\ i).truncLen \leq (open\ j).truncLen \land
        ( (open\ i).truncLen \geq 1 \implies (open\ i).seg1 = (open\ j).seg1 ) \land
        ( (open\ i).truncLen \geq 2 \implies (open\ i).seg2 = (open\ j).seg2 ) ) ) )
\end{schema}

Reading the three disjuncts against \field{entryPairConflicts}
(conflict.go:186--217): the first is \texttt{ws-dir}$\times$\texttt{ws-dir}
or \texttt{ws-file}$\times$\texttt{ws-file}, dispatched to
\field{pathsOverlap}'s truncate-and-compare rule
(conflict.go:263--289) --- a root glob on either side overlaps
unconditionally (``\emph{an entry whose FIRST segment globs claims
from the root}'', conflict.go:303--305), otherwise the truncated
literal prefixes must agree position-for-position up to the shorter
one's length. The second and third are the two orderings of
\texttt{ws-file}$\times$\texttt{ws-dir} (and, by \field{ENTRY}'s own
invariant forcing \field{isEI = ztrue \implies isDir = ztrue}, also
cover \texttt{ws-file}$\times$\texttt{ei-dir}), dispatched to
\field{pathContainedBy}. The leading \field{isEI}-exclusion is
\texttt{ei-dir}$\times$\texttt{ei-dir}, which
\field{entryPairConflicts} always refuses (conflict.go:196--202):
``\emph{Two missions may extract into the same directory... same-filename
collisions are the leader's responsibility, not admission control's}.''

\subsection{Init}

\begin{schema}{Init}
Store'
\where
open' = \emptyset
\end{schema}

\subsection{AdmitMission}

Models \field{Store.Create}'s admission gate: the new entry may claim
an empty slot only if it overlaps no currently-open mission's entry.

\begin{schema}{AdmitMission}
\Delta Store \\
slot? : 1 \upto maxOpen \\
entry? : ENTRY
\where
slot? \notin \dom open \\
\forall i : \dom open @
  \lnot ( \lnot (entry?.isEI = ztrue \land (open\ i).isEI = ztrue) \land
    ( (entry?.isDir = (open\ i).isDir \land
        entry?.truncLen > 0 \land (open\ i).truncLen > 0 \land
        ( entry?.rootGlob = ztrue \lor (open\ i).rootGlob = ztrue
          \lor ( entry?.seg1 = (open\ i).seg1 \land
                ( entry?.truncLen = 1 \lor (open\ i).truncLen = 1
                  \lor entry?.seg2 = (open\ i).seg2 ) ) ) )
      \lor (entry?.isDir = zfalse \land (open\ i).isDir = ztrue \land
        (open\ i).truncLen \leq entry?.truncLen \land
        ( (open\ i).truncLen \geq 1 \implies (open\ i).seg1 = entry?.seg1 ) \land
        ( (open\ i).truncLen \geq 2 \implies (open\ i).seg2 = entry?.seg2 ) )
      \lor ((open\ i).isDir = zfalse \land entry?.isDir = ztrue \land
        entry?.truncLen \leq (open\ i).truncLen \land
        ( entry?.truncLen \geq 1 \implies entry?.seg1 = (open\ i).seg1 ) \land
        ( entry?.truncLen \geq 2 \implies entry?.seg2 = (open\ i).seg2 ) ) ) ) \\
open' = open \oplus \{ slot? \mapsto entry? \}
\end{schema}

The guard is syntactically identical to \field{Store}'s own
invariant with \field{entry?} standing in for \field{open\ i} and
\field{open\ j} ranging only over the currently-open set: exactly the
shape a calculated precondition should have, since the only way
\field{open}$'$ can violate \field{Store}$'$'s invariant is for the
newly-admitted entry to overlap something already there.

\subsection{ReleaseMission}

Not one of \field{findWriteSetConflicts}'s own concerns --- no
function in \fpath{internal/mission} is named this --- but every real
mission eventually leaves the open set (\field{Store.Close},
\field{Store.Abandon}, spec-mission-lifecycle.tex
\S\ref{sec:opclose}--\ref{sec:opabandon}), and without some way to
free a slot the admission model could only ever fill up, never
demonstrate the invariant surviving a longer, mixed sequence of
admits and departures. Removing an entry from \field{open} can only
shrink the set the invariant quantifies over, so it can never
introduce a violation.

\begin{schema}{ReleaseMission}
\Delta Store \\
slot? : 1 \upto maxOpen
\where
slot? \in \dom open \\
open' = \{ slot? \} \ndres open
\end{schema}

% ===================================================================
\section{Theorems}
\label{sec:theorems}
% ===================================================================

\begin{center}
\begin{tabularx}{\textwidth}{@{}lX@{}}
\toprule
\textbf{Theorem} & \textbf{Justification} \\
\midrule
\field{SegmentationUnchanged} &
  \field{CanonicalPathBuggy} and \field{CanonicalPathFixed} compute
  the identical \field{segs} = \field{squash(raw \comp unwrap)} from
  the same \field{raw}; the fix adds a field, it does not change how
  segments are derived. Confirmed by inspection of both schemas
  (\S\ref{sec:canonical}); not separately animated, since the two
  schemas share the defining conjunct verbatim. \\
\field{AbsolutePreserved} &
  Criterion (a). \field{CanonicalPathFixed}'s \field{absolute} field
  is exactly the leading-\field{rsEmpty} test applied to \field{raw},
  by its own defining predicate, and
  \field{PathContainedByFixed} consults it
  (\field{absolute = zfalse}) before consulting containment at all.
  Exhibited as a reachable-deadlock counter-example against the
  \emph{buggy} pairing, and as a clean run against the fixed pairing,
  in \S\ref{sec:validation-vaib}. \\
\field{ContainmentIsSemantic} &
  Criterion (b). \field{PathContainedByBuggy}/\field{Fixed} both
  consult \field{ENTRY}'s \field{seg1}/\field{seg2}/\field{truncLen}
  fields, which are already glob-resolved
  (\field{ENTRY}'s own predicate, \S\ref{sec:state}), never
  \field{pattern} compared as literal text against \field{segs}. The
  contrasting, string-literal reading --- what mission lint's
  coverage heuristic does instead --- is exhibited as a reachable
  deadlock in \S\ref{sec:validation-8ady}. \\
\field{WriteSetNonOverlap} &
  Criterion (c). \field{Store}'s own invariant states it;
  \field{AdmitMission}'s guard is that invariant's calculated
  precondition (\S\ref{sec:admission}); \field{ReleaseMission} can
  only shrink \field{dom\ open}, so it preserves any invariant that
  universally quantifies over that domain. Additionally exercised by
  ProB (\S\ref{sec:validation-admission}): certified complete in
  isolation (\textbf{7266 states, 35915 transitions}), and confirmed
  again, uncertified but reaching the identical state count, against
  the committed file as a whole (\textbf{7266 states, 43180
  transitions}); zero violations in either run. \\
\bottomrule
\end{tabularx}
\end{center}

None of the four are discharged by proof --- as
spec-mission-lifecycle.tex's own Theorems section puts it, ``\fuzz{}
type-checks schemas; it does not discharge proof obligations.'' The
last is the only one whose \emph{own} state space was exhaustively
searched by ProB rather than inspected by hand; the first three are
inspection claims about schemas checked individually, each backed by
a separate, decoupled reachability demonstration in
\S\ref{sec:validation}.

% ===================================================================
\section{Validation}
\label{sec:validation}
% ===================================================================

Three separate probcli runs back the claims above. Each of
\S\ref{sec:validation-vaib} and \S\ref{sec:validation-8ady} follows
spec-mission-lifecycle.tex's own decoupled-demonstration technique
(its \S~``A decoupled demonstration''): a small, self-contained Z
document, checked and animated independently of this file's own
namespace, built so that the safety property being tested is
\emph{not} baked into the type the buggy operation is checked
against. That errata is worth restating in full, because it is the
reason these demonstrations are built the way they are and not as
extra schemas typed against \field{ENTRY}/\field{Store} above: a
schema written as $\Delta S \land P$ silently carries \field{S}'s
\emph{own} invariant on its primed components regardless of what $P$
says, so if the safety property already lives in \field{S}'s
predicate, a buggy $P$ does not violate it --- it has no valid
after-state at all, and ProB correctly (and uselessly) reports
nothing wrong. \field{Store}'s own invariant is exactly this kind of
property, so neither demonstration below is typed against
\field{ENTRY} or \field{Store}; each builds its own minimal,
invariant-free state type and converts the safety property into a
reachability question via the same \field{Idle} self-loop convention
spec-mission-lifecycle.tex uses: enabled everywhere except the one
state the demonstration exists to find, so a buggy operation reaching
that state shows up as a deadlock, and a fixed operation's clean run
shows the state is never reached at all.

The two demonstration files are scratch --- built under
\fpath{.tmp/}, checked and animated, and not committed, exactly as
spec-mission-lifecycle.tex's own two decoupled documents were not
committed alongside it. Their content is reproduced here verbatim
for the record.

\subsection{ethos-vaib: absolute paths swallowed by a relative entry}
\label{sec:validation-vaib}

Shared state schema, carrying \field{CanonicalPath}'s own logic
(\S\ref{sec:canonical}) plus a fixed write-set entry standing in for
\texttt{docs/**}, and a boolean recording whether the buggy
containment check (\S\ref{sec:containment}, ignoring \field{absolute}
entirely) admits the resulting path:

\begin{verbatim}
SEG ::= segDocs | segXmd
RAWSEG ::= rsSeg <SEG> | rsEmpty | rsDot
PATSEG ::= litSeg <SEG> | wildRestSeg
ZBOOL ::= ztrue | zfalse

maxSegs : NN;  maxSegs = 2
unwrap : RAWSEG -+> SEG
  dom unwrap = ran rsSeg              -- undefined on rsEmpty, rsDot
  forall s : SEG @ unwrap(rsSeg(s)) = s
witnessPattern : seq PATSEG
  witnessPattern = <litSeg(segDocs), wildRestSeg>

FileRaw
  raw : seq RAWSEG
  segs : seq SEG
  absolute : ZBOOL
  covered : ZBOOL       -- witnessPattern contains segs, glob-aware
where
  #raw <= maxSegs
  segs = squash(raw ; unwrap)
  absolute = (if #raw > 0 and raw(1) = rsEmpty then ztrue else zfalse)
  covered = ztrue <=> (witnessPattern contains segs, per Contains)

Init
  FileRaw'

Idle                                     -- buggy pairing
  Xi FileRaw
where
  not (absolute = ztrue and covered = ztrue)
\end{verbatim}

The \field{dom unwrap = ran rsSeg} line matters and was missing from
an earlier round of this scratch file: without it, \field{unwrap}'s
axiom pins down its value only on \field{rsSeg}-wrapped inputs and
leaves whether it is \emph{also} defined on \field{rsEmpty}/\field{rsDot}
completely open, so \texttt{SETUP\_CONSTANTS} is free to manufacture
extra, spurious witnesses for it --- four of them, at the setsize this
demonstration otherwise uses, inflating every count below without
changing which state deadlocks. The line closes that gap by stating
\field{unwrap}'s domain exactly.

\texttt{/z-spec:check} passes on the scratch file with this content.
\texttt{/z-spec:model\_check} still reports a deadlock, now with a
single \texttt{SETUP\_CONSTANTS} witness rather than four: \textbf{1
state, 5 transitions} at setsize~2, \textbf{2 states, 6 transitions}
at setsize~3. A deadlock search halts at the first violation it
finds, so neither figure is the full space's size --- only the
deadlock's existence is a robust claim here, exactly as jra's review
of the previous round noted; what the fix changes is that the count is
no longer inflated by an under-constrained function with no bearing on
the bug. The counter-example itself is unchanged:

\begin{verbatim}
1  SETUP_CONSTANTS
2  INITIALISATION      (raw := <rsEmpty, rsSeg(segDocs)>)
                        -- segs = <segDocs>, absolute = ztrue,
                        --  covered = ztrue (witnessPattern's
                        --  literal prefix <docs> matches segs'
                        --  one element; ** swallows the rest)
   -- deadlock: no operation enabled --
   absolute = ztrue, covered = ztrue
\end{verbatim}

The violating valuation is not printed verbatim by the tool's
structured summary (it reports the deadlock and the operation
sequence, not the satisfying assignment); it is given here because it
follows directly from \field{FileRaw}'s own defining equations ---
\field{raw} = $\langle \field{rsEmpty}, \field{rsSeg(segDocs)} \rangle$
is the \emph{only} length-$\leq 2$ value whose \field{absolute} field
is \field{ztrue} and whose \field{segs} is contained by
\field{witnessPattern}, so it is the state the search finds regardless
of setsize. This is \field{ethos-vaib} exactly: an absolute path ---
\texttt{"/docs"}, in the ticket's own terms --- treated by the buggy
containment check as though it were the repo-relative \texttt{docs},
and admitted by the glob entry \texttt{docs/**} that was only ever
meant to authorize writes under the repo's own \texttt{docs/}
directory.

The corrected pairing changes only the \field{Idle} guard, to require
\field{fixedAdmit} --- containment additionally gated on
\field{absolute = zfalse}, matching \field{PathContainedByFixed}
exactly:

\begin{verbatim}
Idle                                     -- corrected pairing
  Xi FileRaw
where
  not (absolute = ztrue and fixedAdmit = ztrue)
  -- fixedAdmit = ztrue <=> (absolute = zfalse and covered = ztrue)
\end{verbatim}

\texttt{/z-spec:check} passes; \texttt{/z-spec:model\_check} explores
\textbf{5 states, 9 transitions} at both setsize~2 and setsize~3 ---
stable, unlike the buggy pairing above, because a clean run keeps
searching until the space is exhausted rather than stopping at a
first violation --- and reports \textbf{no} deadlock and \textbf{no}
violation (probcli does not mark it certified complete at this size,
the same caveat \S\ref{sec:validation-admission} records for the
full admission model). The state \field{Idle} excludes is never
reached, because \field{fixedAdmit} can never be \field{ztrue} at the
same time as \field{absolute = ztrue} --- not because a guard happens
to check for it at runtime, but because the two conjuncts of
\field{fixedAdmit}'s own definition are contradictory whenever
\field{absolute = ztrue}.

\subsection{ethos-8ady: a glob compared as literal text}
\label{sec:validation-8ady}

Mission lint's coverage heuristic
(\fpath{internal/mission/lint.go}:191--197) does not compare every
write-set entry against the candidate file: the loop's own guard,
line~193, is $\field{strings.HasSuffix(w, "/")} \land
\field{strings.HasPrefix(cf, CanonicalPath(w)+"/")}$, and Go's
\field{HasSuffix} short-circuits the whole comparison for an entry
that does not end in \texttt{"/"}. \texttt{docs/**} does not ---
it ends in \texttt{*} --- so for exactly the entry this demonstration
uses, the \field{HasPrefix} half is never evaluated at all: the entry
is silently skipped, not compared and found wanting.
\field{witnessEntryIsDir} below is that guard made explicit, fixed to
\field{zfalse} for this witness; \field{Stringify} is the round trip
\field{CanonicalPath} would still perform on \field{witnessEntry}'s
segments \emph{if} the guard passed, kept here so \field{literalCovered}
states the real predicate's full shape rather than silently dropping
the half that never fires for this witness.

\begin{verbatim}
SEG ::= segDocs | segFile | segStarStar
PATSEG ::= litSeg <SEG> | wildRestSeg
ZBOOL ::= ztrue | zfalse

maxSegs : NN;  maxSegs = 2
Stringify : PATSEG -> SEG
  Stringify(wildRestSeg) = segStarStar
  forall s : SEG @ Stringify(litSeg(s)) = s
witnessEntry : seq PATSEG
  witnessEntry = <litSeg(segDocs), wildRestSeg>
witnessEntryIsDir : ZBOOL             -- strings.HasSuffix("docs/**", "/")
  witnessEntryIsDir = zfalse          --   = false: no trailing slash

FileSegs
  segs : seq SEG
  semanticCovered : ZBOOL    -- witnessEntry contains segs, glob-aware
  literalCovered : ZBOOL     -- lint.go:193's guarded HasPrefix check
where
  #segs <= maxSegs
  semanticCovered = ztrue <=> (witnessEntry contains segs, per Contains)
  literalCovered = ztrue <=>
    witnessEntryIsDir = ztrue and
    #witnessEntry < #segs and
    (forall i : 1..#witnessEntry @ Stringify(witnessEntry(i)) = segs(i))

Init
  FileSegs'

Idle                                     -- buggy pairing
  Xi FileSegs
where
  not (semanticCovered = ztrue and literalCovered = zfalse)
\end{verbatim}

\texttt{/z-spec:check} passes. \texttt{/z-spec:model\_check} explores
\textbf{1 state, 5 transitions} and reports a deadlock immediately,
at \field{INITIALISATION}. \field{literalCovered} is \field{zfalse}
for \emph{every} value of \field{segs}, unconditionally, because
\field{witnessEntryIsDir = zfalse} already fails the guard's leading
conjunct --- exactly \field{HasSuffix("docs/**", "/")} returning
\field{false} in the real code, before \field{HasPrefix} is ever
reached. \field{semanticCovered}, by contrast, is \field{ztrue} for
$\field{segs} = \langle \field{segDocs}, \field{segFile} \rangle$
(the file \texttt{docs/audited-delegation.md}, in the same terms
\field{pathContainedBy}'s own doc comment uses, conflict.go:339--340)
--- the \texttt{**} swallows the second segment. This is
\field{ethos-8ady}: mission lint reports
\texttt{docs/audited-delegation.md} as \emph{not} covered by
\texttt{docs/**}, when it plainly is, because \field{docs/**} is
never even a candidate the coverage loop considers.

The corrected pairing routes the heuristic through the same
\field{Contains} logic \field{PathContainedByFixed} uses
(matching the real fix bd-recommends: ``\emph{route lint.go coverage
through the same mission.PathContainedBy matcher}''), which collapses
the guard to a tautology:

\begin{verbatim}
Idle                                     -- corrected pairing
  Xi FileSegs
where
  not (semanticCovered = ztrue and semanticCovered = zfalse)
\end{verbatim}

\texttt{/z-spec:check} passes; \texttt{/z-spec:model\_check} explores
\textbf{5 states, 9 transitions} and reports \textbf{no} deadlock and
\textbf{no} violation --- the two notions of ``covered'' cannot
disagree once lint asks the same question \field{PathContainedByFixed}
does.

\subsection{The write-set admission model}
\label{sec:validation-admission}

Unlike the two demonstrations above, \field{Store} and its operations
(\S\ref{sec:admission}) are checked \emph{in this document's own
namespace} --- the errata's trap does not apply here, because
\field{Store}'s invariant is the property being \emph{confirmed}, not
a buggy alternative being checked against a type that already
assumes it. \texttt{/z-spec:check} reports 0 errors on the committed
file.

Run against the committed file as a whole, \texttt{/z-spec:model\_check}
necessarily also enumerates \field{PathContainedByBuggy}'s and
\field{PathContainedByFixed}'s own domains alongside \field{Store}'s
(\S\ref{sec:modellingchoice} explains why these coexist safely in one
namespace despite each being $\seq X$-typed: they are schemas, not
free functions, so \texttt{SETUP\_CONSTANTS} never has to solve for
them as constants). That combined run explores \textbf{7266 states,
43180 transitions} --- the identical state count the isolated run
below finds, with more transitions because
\field{PathContainedByBuggy}/\field{Fixed} fire as additional
enumerated schemas without contributing any new \field{Store} state
--- with zero invariant violations and zero deadlocks, though probcli
does not report it as certified complete at that size (raising the
operation budget to 500000 changed nothing). Isolating
\field{Init}/\field{AdmitMission}/\field{ReleaseMission} alone ---
the same three schemas, in a scratch file carrying no other
declaration --- \texttt{/z-spec:model\_check} explores \textbf{7266
states, 35915 transitions} and \emph{is} reported certified complete,
with the same result: zero invariant violations, zero deadlocks. Both
runs agree on the property that matters, and on the size of the
state space it holds across: \field{Store}'s pairwise non-overlap
invariant holds everywhere \field{AdmitMission} and
\field{ReleaseMission} can reach.

\subsection{Tool run}
\label{sec:probrun}

\begin{center}
\begin{tabularx}{\textwidth}{@{}lccl@{}}
\toprule
\textbf{Model} & \textbf{States} & \textbf{Transitions} & \textbf{Result} \\
\midrule
Committed file as a whole (\field{Store} +
  \field{PathContainedByBuggy}/\field{Fixed}) & 7266 & 43180 & 0
  violations, not certified complete \\
\field{Init}/\field{AdmitMission}/\field{ReleaseMission} in
  isolation & 7266 & 35915 & certified complete, 0 violations \\
Decoupled, \field{ethos-vaib} buggy pairing (setsize 2 / setsize 3) &
  1 / 2 & 5 / 6 & deadlock at
  $\field{absolute} = \field{ztrue} \land \field{covered} = \field{ztrue}$;
  a deadlock search halts at the first violation, so neither count is
  the full space \\
Decoupled, \field{ethos-vaib} corrected pairing & 5 & 9 & clean,
  identical at setsize 2 and 3 \\
Decoupled, \field{ethos-8ady} buggy pairing & 1 & 5 & deadlock at
  \field{INITIALISATION} \\
Decoupled, \field{ethos-8ady} corrected pairing & 5 & 9 & clean \\
\bottomrule
\end{tabularx}
\end{center}

fuzz's \texttt{/z-spec:check} type-checks this document with 0 errors.
Every scratch file above was saved with a \texttt{.tex} extension
before being checked or animated, for the reason
spec-mission-lifecycle.tex's own Tool-run section already recorded:
probcli silently reports 0 states and 0 transitions against any other
extension.

% ===================================================================
\section{Traceability}
\label{sec:trace}
% ===================================================================

\begin{center}
\begin{tabularx}{\textwidth}{@{}llX@{}}
\toprule
\textbf{Z construct} & \textbf{Go source} & \textbf{Note} \\
\midrule
\field{RAWSEG} & \field{splitSegments}'s input shape
  (conflict.go:475--497) & the pre-filter split tokens; not a real Go
  type \\
\field{unwrap} & the \field{s == ""}/\field{s == "."} filter inside
  \field{splitSegments} (conflict.go:487--491) & \\
\field{CanonicalPathBuggy} & \field{splitSegments} +
  \field{CanonicalPath} (conflict.go:475--497, 518--524) & today's
  real behavior \\
\field{CanonicalPathFixed} & the fix bd ticket \field{ethos-vaib}
  prescribes & not present in current code \\
\field{PATSEG}, \field{ENTRY.pattern} & entry segments after
  splitting, glob-classified & \field{segmentMatches}
  (conflict.go:433--442) does the classification per-segment; not
  separately modelled (\S\ref{sec:nongoals}) \\
\field{ENTRY.isDir} & \field{isDirEntry} (conflict.go:219--226) & \\
\field{ENTRY.isEI = ztrue \implies isDir = ztrue} & validate.go rule
  17 (validate.go:79, 216--226) & \\
\field{ENTRY.rootGlob}, \field{truncLen}, \field{seg1}, \field{seg2} &
  \field{overlapSegments} (conflict.go:291--314) & animatability
  device, \S\ref{sec:state} \\
\field{PathContainedByBuggy} & \field{pathContainedBy} +
  \field{segmentsContain} (conflict.go:360--379, 394--419) & \\
\field{PathContainedByFixed} & the fix bd ticket \field{ethos-vaib}
  prescribes, applied to \field{pathContainedBy} & not present in
  current code \\
\field{Store} & \field{Store.Create}'s existing-open-missions set,
  as consulted by \field{findWriteSetConflicts}
  (conflict.go:70--120) & \\
\field{Store}'s invariant, \field{AdmitMission}'s guard &
  \field{anyEntryConflicts}/\field{entryPairConflicts}
  (conflict.go:154--217), \field{pathsOverlap}
  (conflict.go:263--289) & \\
\field{ReleaseMission} & \field{Store.Close}/\field{Store.Abandon}
  freeing a slot & no single named Go function; see
  spec-mission-lifecycle.tex \field{Close}/\field{Abandon} \\
Literal-string coverage contrast, \field{witnessEntryIsDir} &
  \fpath{internal/mission/lint.go}:193's guard, exactly
  (\S\ref{sec:validation-8ady}) & \field{ethos-8ady}; a direct
  transcription of the guard, not a representative abstraction --- see
  \S\ref{sec:validation-8ady} for why \texttt{docs/**} never reaches
  the \field{HasPrefix} half at all \\
\bottomrule
\end{tabularx}
\end{center}

\end{document}
